\subsection{Phase geometry and measured profile responses}
\label{sec:three-band-chain}

For $s>1$, the installed frequency $\omega'_j=\omega_j s^{-m_j}$ maps separation $\Delta$ to native effective separation $\Delta s^{-m_j}$. At $\Delta=s\Ltr$, containment within the unwrapped training interval requires $m_j\ge1$. This is a per-pair geometric condition: distinct $m_j$ generally map to distinct native separations, rather than one shared training position. The phase sensitivity at the native table is
\[
\left.\frac{\partial(\Delta\omega_j s^{-m_j})}{\partial m_j}\right|_{m_j=0}
=-\Delta\omega_j\ln s.
\]
It quantifies local phase change; the finite drift at $m_j=0$ is zero.

Proposition~\ref{prop:collapse} characterizes positional directions, while the learned content coefficients can differ across pairs. We assess deployment profiles through complete-generation task scores.

For twelve OLMo-$1$B profiles on $350$ shared $16$K RULER inputs, a count of slots whose test-distance turns lie in $[0.25,16]$ separates three response groups. Appendix~\ref{sec:profile-diagnostic-audit} defines the count, tabulates every arm, and reports its reconstruction provenance. The count is an exploratory summary of this panel: it separates between-group scores while leaving within-group rankings unresolved.
